% !TEX root = ../../main.tex
% =====================================================================
% Chapter: Arnold residues and the typed scalar ledger
%
% The chapter records two comparison mechanisms:
% (G2) a finite-window KZ--Arnold bar-superconnection comparison;
% (G3) a typed ledger separating central, modular, ghost, lattice,
%      and Borcherds scalars.
%
% Conventions (visible at every relevant identity):
% - KZ residue in the affine conformal-block window:
%   Omega / (k + h^v).
% - kappa values: read from chapters/examples/landscape_census.tex
% (source every formula from the manuscript or a primary reference)
% - Vol III kappa: subscripted (kappa_ch / kappa_BKM / kappa_cat /
% kappa_fiber); cross-volume identity refers to the bridge
% mediated by Phi.
% - Bar complex: B^ord(A) = T^c(s^{-1} A-bar).
% =====================================================================

\providecommand{\fg}{\mathfrak{g}}
\providecommand{\fh}{\mathfrak{h}}
\providecommand{\hv}{h^{\vee}}
\providecommand{\Conf}{C}
\providecommand{\Confnord}{\Conf_n^{\mathrm{ord}}}
\providecommand{\BarOrd}{B^{\mathrm{ord}}}
\providecommand{\BarSym}{B^{\Sigma}}
\providecommand{\Tcoalg}{T^{c}}
\providecommand{\Sym}{\mathrm{Sym}}
\providecommand{\Op}{\mathrm{Op}}
\providecommand{\End}{\mathrm{End}}
\providecommand{\Hom}{\mathrm{Hom}}
\providecommand{\Res}{\operatorname{Res}}
\providecommand{\Ran}{\mathrm{Ran}}
\providecommand{\BRST}{\mathrm{BRST}}
\providecommand{\cghost}{c_{\mathrm{ghost}}}
\providecommand{\bc}{\mathrm{bc}}
\providecommand{\Sug}{\mathrm{Sug}}
\providecommand{\KZ}{\mathrm{KZ}}
\providecommand{\KZB}{\mathrm{KZB}}
\providecommand{\Arn}{\mathrm{Arnold}}
\providecommand{\ConnConf}{\mathsf{ConnConf}}
\providecommand{\ChirAlg}{\mathsf{ChirAlg}}
\providecommand{\BRSTGaugedChirAlg}{\mathsf{BRSTGaugedChirAlg}}
\providecommand{\Cone}{\operatorname{cone}}
\providecommand{\dlog}{d\log}
\providecommand{\Phihol}{\Phi_{\mathrm{hol}}}
\providecommand{\PhiCY}{\Phi}
\providecommand{\SpCh}{\mathrm{Sp}^{\mathrm{ch}}}
\providecommand{\PhiFA}{\Phi^{\mathrm{FA}}}
\providecommand{\dbar}{d_{\mathrm{bar}}}
\providecommand{\Hch}{H^{\bullet}_{\mathrm{ch}}}
\providecommand{\IdC}{\mathrm{Id}_{\mathrm{ConnConf}}}

\chapter{Arnold Residues and Scalar Invariants}
\label{chap:climax-platonic}\label{ch:chiral-climax-platonic}
\index{climax theorem!typed comparison surface|textbf}
\index{Arnold connection!pullback identity|textbf}
\index{scalar invariants!central, modular, ghost, lattice|textbf}
\index{KZ functor!chiral algebras|textbf}

\ChapterIntro{Theorems A, B, C, D, H traverse Volume~I in five
categorical directions: bar--cobar adjunction, completed inversion,
conductor complementarity, genus-$g$ integration, and chiral
Hochschild cohomology. Their proofs live in their native chapters.
This chapter supplies a common notation for two typed statements:
\[
 \nabla^{B,0}_{A,V,n}
 =d_A+d_{\mathrm{dR}}
 -\sum_{i<j}\widetilde\rho_{A,V,n}(t_{ij})\,\eta_{ij}
 =\KZ_{A,V,n}^{*}(\nabla_{\Arn}),
 \qquad
 \bigl(K^c,\,K^\kappa,\,c_{\mathrm{gh}},\,K^{\mathrm{lat}},
 \kappa_{\mathrm{BKM}}\bigr).
\]
The first statement is conditional on the logarithmic-residue package
$H_{\mathrm{KZ}}$: a finite bar window, logarithmic collision
operators, the infinitesimal braid relations, and the
Fulton--MacPherson comparison.  It identifies the finite-window
ordered genus-$0$ bar superconnection with the pullback of Arnold's
$1969$ operator-valued KZ connection; the bar differential is its
Fulton--MacPherson residue realization.  Drinfeld--Kohno supplies the
standard affine example.

The second display separates the scalar invariants by ambient.  It keeps the
central-charge conductor $K^c$, the modular conductor $K^\kappa$, the
central charge $c_{\mathrm{gh}}$ of a specified BRST ghost complex,
the lattice number $K^{\mathrm{lat}}=2c_+$, and the Borcherds weight
$\kappa_{\mathrm{BKM}}$ in their native ambients.  The scalar values
computed from the named companion formulas are
$\{0,13,250/3\}$.  Bershadsky--Polyakov contributes the exact
central-charge identity $K^c_{\mathrm{BP}}=50$; its strong generators
$J,G^+,G^-,T$ are even, and its modular conductor remains an open
genus-one computation.  The Mukai lattice has signature $(4,20)$ and
therefore $K^{\mathrm{lat}}=2c_+=8$; the chiral equation
$K^\kappa=8$ is conjectural under
$H_{\mathsf B}=(H_{\mathrm{chart}},H_{\mathrm{KD}},H_{\mathrm{scalar}},
H_{\mathrm{mod}},H_{\mathrm{quantum}})$.  These distinctions govern
the five theorem pointers and the Vol~II/III comparison remarks below.}

\section{Setup: the Arnold connection, the KZ functor, BRST resolution}
\label{sec:climax-platonic-setup}

Four objects enter the comparison: the ordered configuration space
$\Confnord(\Conf)$ with Arnold's $1$-forms; the universal
infinitesimal braid Lie algebra $\mathfrak{t}_n$; the chiral bar
coalgebra $\BarOrd(A)$; and the ghost charge of a quasi-free BRST
resolution. Each is recalled in the precise form the two comparison
surfaces require.

\subsection{The Arnold connection}
\label{ssec:climax-platonic-arnold}

\begin{lemma}[Arnold three-term form identity; \ClaimStatusProvedElsewhere]
\label{lem:arnold-three-term-forms-platonic}
For the Arnold part of the climax, $\Conf$ denotes an affine
genus-$0$ coordinate screen, equivalently a formal disk inside
$\mathbb{P}^1$, with coordinate~$z$. Let
$\Confnord(\Conf) = \{(z_1,\dots,z_n) \in \Conf^{n}: z_i \neq z_j\}$
denote the open configuration space of $n$ ordered distinct points.
The Arnold $1$-forms on this screen are
\begin{equation}
\label{eq:arnold-eta}
 \eta_{ij} \;=\; \dlog(z_i - z_j)
 \;\in\; \Omega^{1}\bigl(\Confnord(\Conf)\bigr),
 \qquad 1 \leq i \neq j \leq n.
\end{equation}
This is an Arnold-local affine/formal representative of the
logarithmic normal form along the diagonal. On a curve it must be
glued by the coordinate-change cocycle; at positive genus it is
replaced by the KZB/theta or prime-form representative.
Arnold's three-term identity reads
\begin{equation}
\label{eq:arnold-three-term}
 \eta_{ij} \wedge \eta_{jk}
 \;+\; \eta_{jk} \wedge \eta_{ki}
 \;+\; \eta_{ki} \wedge \eta_{ij}
 \;=\; 0
 \quad\text{in }\Omega^{2}\bigl(\Conf_3^{\mathrm{ord}}(\Conf)\bigr),
\end{equation}
for every triple of distinct indices.
\end{lemma}

\begin{proof}[Attribution]
This is the logarithmic-form relation in Arnold's braid-arrangement
calculation; the cohomological version and the Arnold--Orlik--Solomon
attribution are recalled immediately below.
\end{proof}

\begin{theorem}[Arnold 1969]
\label{thm:arnold-cohomology-platonic}
\ClaimStatusProvedElsewhere
For $\Conf=\mathbb{A}^1$ the de Rham cohomology ring
$H^{*}(\Confnord(\mathbb{A}^1);\mathbb{C})$ is the
graded-commutative algebra generated by the classes $[\eta_{ij}]$,
subject to $\eta_{ij}=\eta_{ji}$, $\eta_{ij}^{2}=0$, and the
three-term identity~\eqref{eq:arnold-three-term}. On
$\mathbb{P}^1$ the descended pairwise classes also satisfy the
projective linear relations
$\sum_{j\neq i}\bar\eta_{ij}=0$; for positive-genus curves the
Totaro--K\v{r}\'{\i}\v{z}--Bezrukavnikov model adds the surface
classes, mixed period relations, and differential
$d\eta_{ij}=\Delta_{ij}$.
\end{theorem}

\begin{remark}[Attribution]
\label{rem:arnold-attribution-platonic}
Theorem~\ref{thm:arnold-cohomology-platonic} is the Arnold--Orlik--Solomon
description of the cohomology of complements of hyperplane arrangements
applied to the braid arrangement on $\mathbb{A}^1$; see Arnold
(\textit{Mat.\ Zametki}~$5$, $1969$) and Orlik--Solomon
(\textit{Inventiones}~$56$, $1980$). The projective and positive-genus
corrections are the genus split of
Proposition~\ref{prop:arnold-higher-genus}.
\end{remark}

\begin{definition}[Infinitesimal braid Lie algebra; \ClaimStatusDefinitional]
\label{def:infinitesimal-braid-lie-algebra-platonic}
The infinitesimal braid Lie algebra $\mathfrak{t}_n$ is the free
$\mathbb{Z}$-graded Lie algebra on generators
$t_{ij} = t_{ji}$ for $1 \leq i \neq j \leq n$ modulo the
infinitesimal braid relations
\begin{equation}
\label{eq:inf-braid-relations}
 [t_{ij}, t_{kl}] = 0 \quad \text{when } \{i,j\} \cap \{k,l\} = \emptyset,
 \qquad
 [t_{ij},\, t_{ik} + t_{jk}] = 0 \quad \text{for distinct } i,j,k.
\end{equation}
\end{definition}

\begin{definition}[Arnold's universal KZ connection; \ClaimStatusDefinitional]
\label{def:arnold-universal-kz-connection-platonic}
\emph{Arnold's universal KZ connection} on the trivial
$U(\mathfrak{t}_n)$\nobreakdash-bundle over $\Confnord(\Conf)$ is
\begin{equation}
\label{eq:nabla-arnold}
 \nabla_{\Arn}
 \;=\; d \;-\; \sum_{1 \leq i < j \leq n} t_{ij}\,\eta_{ij},
\end{equation}
with $\eta_{ij}$ the Arnold $1$-form~\eqref{eq:arnold-eta} and
$t_{ij} \in \mathfrak{t}_n$ the universal infinitesimal braid generator.
\end{definition}

\begin{proposition}[Arnold flatness]
\label{prop:arnold-flatness-platonic}
\ClaimStatusProvedElsewhere
The connection $\nabla_{\Arn}$ is flat: $\nabla_{\Arn}^{2} = 0$. Flatness
is equivalent to the simultaneous vanishing of the three-term
identity~\eqref{eq:arnold-three-term} and the infinitesimal braid
relations~\eqref{eq:inf-braid-relations}.
\end{proposition}

\begin{proof}[Attribution]
The flatness $\nabla_{\Arn}^{2} = 0$ unpacks to
$\sum [t_{ij}, t_{kl}]\,\eta_{ij}\wedge\eta_{kl} = 0$, and the
hypotheses on $t_{ij}$ together with~\eqref{eq:arnold-three-term}
make every wedge sum vanish. Reference: Kohno (\textit{Annales
Inst.\ Fourier}~$37$, $1987$), Drinfel'd (\textit{Algebra i
Analiz}~$2$, $1989$).
\end{proof}

\subsection{The chiral bar complex $\BarOrd(A)$}
\label{ssec:climax-platonic-bar}

Let $A$ be an $E_\infty$\nobreakdash-chiral algebra on $\Conf$ in
the sense of Beilinson--Drinfeld: an augmented unital factorisation
$D$\nobreakdash-module on $\Conf$ with chiral multiplication
$\mu \colon j_!(A \boxtimes A)|_{\Conf^{2} \setminus \Delta} \to
\Delta_{!} A$. Write $\bar A = \ker\epsilon$ for the augmentation ideal.

\begin{definition}[Ordered chiral bar complex]
\label{def:bar-ord-platonic}
The \emph{ordered chiral bar coalgebra} of $A$ is the cofree conilpotent
graded coassociative coalgebra
\[
 \BarOrd(A)
 \;=\; \Tcoalg\bigl(s^{-1}\bar A\bigr)
 \;=\; \bigoplus_{n \geq 0}\bigl(s^{-1}\bar A\bigr)^{\otimes n},
\]
with deconcatenation coproduct, equipped with the bar differential
$\dbar = d_0 + d_1$ assembled from the augmentation $d_0$ and the
chiral multiplication $d_1$ via the standard Stasheff--Quillen
twisting morphism (Vol~I Chapter~\ref{chap:bar-construction}).
\end{definition}

The bar differential satisfies $\dbar^{2} = 0$.  Its bar--cobar
comparison with $A$ is recorded in
Theorem~\ref{thm:theorem-a-platonic-pointer} below for the bar--cobar
adjoint equivalence on the Koszul locus.

\subsection{The logarithmic KZ comparison datum}
\label{ssec:climax-platonic-kz-functor}

Arnold's connection is universal for representations of the
infinitesimal braid Lie algebra.  A chiral algebra enters this
construction after its collision operators have been shown to define
such a representation.  This order of construction separates the
geometric logarithmic forms from the algebraic Borcherds identities.

\begin{definition}[The logarithmic Arnold arena; \ClaimStatusDefinitional]
\label{def:connconf-platonic}
For fixed $n$, let $\ConnConf^{\mathrm{Arn,log}}_{\Conf,n}$ be the
category whose objects are finite-rank dg vector spaces $(V,d_V)$
equipped with a dg Lie representation
\(
 \rho\colon\mathfrak t_n\to\End(V)
\)
and the associated logarithmic connection
\begin{equation}
\label{eq:nabla-n-form}
 \nabla_{\rho}
 =d_V+d_{\mathrm{dR}}
 -\sum_{1\leq i<j\leq n}\rho(t_{ij})\,\eta_{ij}.
\end{equation}
Morphisms are chain maps intertwining the $\mathfrak t_n$ actions.
The compatible collection over all $n$ is denoted
$\ConnConf^{\mathrm{Arn,log}}_{\Conf}$.
\end{definition}

\begin{remark}[Universal coefficient property; \ClaimStatusProvedHere]
\label{rem:connconf-initial-platonic}
Every representation $\rho\colon\mathfrak t_n\to\End(V)$ extends
uniquely to an algebra morphism
$U(\mathfrak t_n)\to\End(V)$.  Base change of
$\nabla_{\Arn}$ along this morphism is $\nabla_\rho$.  Thus
$U(\mathfrak t_n)$ is the universal coefficient algebra for the
connections in Definition~\ref{def:connconf-platonic}.
\end{remark}

\begin{proof}
This is the universal property of the enveloping algebra.  Substitution
of $t_{ij}\mapsto\rho(t_{ij})$ in~\eqref{eq:nabla-arnold} gives
\eqref{eq:nabla-n-form}; Proposition~\ref{prop:arnold-flatness-platonic}
and the dg-intertwining condition give flatness of the total
superconnection.
\end{proof}

\begin{construction}[Finite-window KZ datum; \ClaimStatusConditional]
\label{constr:kz-functor-platonic}
Let $A$ be an augmented chiral algebra on an affine/formal genus-zero
screen.  A finite window $(V,n)$ of $\BarOrd(A)$ carries the
hypothesis package
\[
H_{\mathrm{KZ}}(A;V,n)
=(H_{\mathrm{fin}},H_{\log},H_{\mathrm{Borch}},H_{\mathrm{FM}}),
\]
with the following content.
\begin{enumerate}[label=\textup{(K\arabic*)}]
\item $H_{\mathrm{fin}}$: $V$ is a finite-rank dg quotient or
subquotient of the length-$\leq n$ ordered bar complex.
\item $H_{\log}$: each pairwise collision has a logarithmic residue
operator $T_{ij}\in\End^0(V)$ commuting with the internal differential.
\item $H_{\mathrm{Borch}}$: the Borcherds mode identities imply
\[
[T_{ij},T_{kl}]=0,
\qquad
[T_{ij},T_{ik}+T_{jk}]=0
\]
for disjoint pairs and distinct triples, respectively.
\item $H_{\mathrm{FM}}$: the Fulton--MacPherson boundary-residue
realization sends the pairwise logarithmic term
$T_{ij}\eta_{ij}$ to the corresponding collision summand of the bar
differential.
\end{enumerate}
Under this package the assignment $t_{ij}\mapsto T_{ij}$ defines
\begin{equation}
\label{eq:r-matrix-from-ope}
 \widetilde\rho_{A,V,n}\colon\mathfrak t_n\longrightarrow\End(V)
\end{equation}
and hence the superconnection
\begin{equation}
\label{eq:nabla-of-A}
 \nabla^{B,0}_{A,V,n}
 =d_A+d_{\mathrm{dR}}
 -\sum_{i<j}\widetilde\rho_{A,V,n}(t_{ij})\eta_{ij}.
\end{equation}
Functoriality holds for morphisms preserving the finite window and the
four pieces of $H_{\mathrm{KZ}}$.
\end{construction}

\begin{computation}[Affine Kac--Moody KZ window; \ClaimStatusProvedElsewhere]
\label{ex:kz-of-vk-g-platonic}
Let $\fg$ be simple, let $k\neq-h^\vee$, and let
$V_1,\ldots,V_n$ be finite-dimensional $\fg$-modules.  The Sugawara
Ward identity for conformal blocks of $V^k(\fg)$ gives
\begin{equation}
\label{eq:r-kz-form-vk-g}
 \nabla_{\mathrm{KZ}}
 =d_{\mathrm{dR}}-\frac{1}{k+h^\vee}
 \sum_{i<j}\Omega_{ij}\,\dlog(z_i-z_j),
\end{equation}
where $\Omega=\sum_a I^a\otimes I_a$ for dual bases of an invariant
form.  The operators
$T_{ij}=\Omega_{ij}/(k+h^\vee)$ satisfy the infinitesimal braid
relations, so this conformal-block window is an object of
$\ConnConf^{\mathrm{Arn,log}}_{\mathbb A^1,n}$.
\end{computation}

\begin{proof}[Attribution]
The current OPE fixes the level and invariant form.  The displayed
connection follows from the Sugawara Ward identity; see
Knizhnik--Zamolodchikov (1984) and Kohno (1987).  Invariance of
$\Omega$ gives the infinitesimal braid relations.  The critical level
$k=-h^\vee$ is precisely the singular Sugawara parameter.
\end{proof}

\subsection{Quasi-free BRST resolution and ghost charge}
\label{ssec:climax-platonic-brst}

The modular characteristic $\kappa(A)$ is defined by the genus-one
curvature of the chiral theory.  A BRST presentation carries a second
datum: the central charge of its ghost fields.  Their comparison is a
theorem upon specification of the presentation, conformal vector, and
normalization map.

Let $\BRSTGaugedChirAlg$ denote the category of pairs $(A, R_{\bullet}
\twoheadrightarrow A)$ consisting of an augmented unital $E_\infty$-chiral
algebra $A$ together with a quasi-free BRST resolution
$R_{\bullet} \to A$ in the homotopical sense: $R_{\bullet}$ is a free
chiral cdga generated in degrees $\leq 0$ by an augmentation-ideal
generator system together with a $bc$-ghost system, the
differential is the BRST charge $Q_{\BRST}$, and the morphism
$R_{\bullet} \to A$ is a quasi-isomorphism.

For each fermionic $bc$ pair of weights $(\lambda,1-\lambda)$, the
central charge is
$c_{\bc}(\lambda)=1-3(2\lambda-1)^2$.  For each bosonic
$\beta\gamma$ pair with the same weights, the central charge is
$c_{\beta\gamma}(\lambda)
= 2(6\lambda^{2} - 6\lambda + 1)$; the sign reverses
relative to $bc$. The total ghost charge of $R_{\bullet}$ is
\begin{equation}
\label{eq:c-ghost-formula}
 \cghost(R_{\bullet})
 \;=\; \sum_{\text{$bc$ pairs}} c_{\bc}(\lambda_{\alpha})
 \;+\; \sum_{\text{$\beta\gamma$ systems}} c_{\beta\gamma}(\lambda_{\beta}).
\end{equation}

\begin{remark}[Statistics and conformal weight]
\label{rem:bc-vs-bg-platonic}
Statistics and conformal weight are separate entries.  The $bc$
system is odd and the $\beta\gamma$ system is even for every value of
$\lambda$.  At $\lambda=2$ their central charges are $-26$ and $+26$,
respectively.  A BRST comparison therefore records both the conformal
vector and the parity of each field.
\end{remark}

\section{Two typed comparison surfaces}
\label{sec:climax-platonic-two-equations}

The operator-valued comparison and the scalar list have different
type signatures.  Their hypotheses are recorded independently.

\subsection{Equation~(G2): bar differential is pullback of Arnold}
\label{ssec:climax-platonic-equation-g2}

Assume $H_{\mathrm{KZ}}(A;V,n)$.  The collision operators define a
representation of $\mathfrak t_n$, and the finite-window bar
superconnection is the corresponding base change of Arnold's
operator-valued connection.  Applying the Fulton--MacPherson residue
realization gives the collision part of~$\dbar$.

\begin{definition}[The pullback identity, structural form; \ClaimStatusDefinitional]
\label{def:pullback-identity-structural-platonic}
The \emph{pullback identity} for the datum $(A;V,n,H_{\mathrm{KZ}})$
on an affine/formal genus-zero collision screen is
\begin{equation}
\label{eq:climax-equation-G2}
\nabla^{B,0}_{A,V,n}
=d_A+d_{\mathrm{dR}}
-\sum_{i<j}\widetilde\rho_{A,V,n}(t_{ij})\,\eta_{ij}
=\KZ_{A,V,n}^{*}(\nabla_{\Arn}),
\end{equation}
where $\KZ_{A,V,n}$ denotes base change along
$U(\mathfrak t_n)\to\End(V)$.  The chain differential $\dbar$ is the
FM-boundary residue realization of this superconnection under
$H_{\mathrm{FM}}$.
\end{definition}

\subsection{Equation~(G3): the scalar type signature}
\label{ssec:climax-platonic-equation-g3}

\begin{definition}[Scalar invariants and their ambients;
\ClaimStatusDefinitional]
\label{def:ghost-identity-structural-platonic}
For a chiral algebra $A$ with a specified companion $A^!$, a specified
BRST complex $R_\bullet$, a lattice $L$ of signature $(c_+,c_-)$,
and a Borcherds product $\Phi$, set
\begin{equation}
\label{eq:climax-equation-G3}
 \mathscr S(A,R_\bullet,L,\Phi)
 :=\left(
 \underbrace{c(A)+c(A^!)}_{K^c(A)},\;
 \underbrace{\kappa(A)+\kappa(A^!)}_{K^\kappa(A)},\;
 \underbrace{c_{\mathrm{gh}}(R_\bullet)}_{\text{BRST presentation}},\;
 \underbrace{2c_+(L)}_{K^{\mathrm{lat}}(L)},\;
 \underbrace{\operatorname{wt}(\Phi)}_{\kappa_{\mathrm{BKM}}(\Phi)}
 \right).
\end{equation}
The first two entries are invariants of a named chiral companion pair;
the third belongs to a chosen BRST presentation; the fourth is lattice
arithmetic; the fifth is an automorphic weight.  A comparison theorem
between two entries specifies the objects, the normalization, and its
hypothesis package.
\end{definition}

For a chosen scalar normalization $\alpha_R\in\mathbb Q$, the BRST
comparison obstruction is
\[
 \Gamma_{\mathrm{gh}}(A,R_\bullet;\alpha_R)
 :=\kappa(A)+\alpha_R c_{\mathrm{gh}}(R_\bullet).
\]
Its vanishing is a family theorem about the displayed presentation.

\section{The comparison theorem}
\label{sec:climax-platonic-climax-theorem}

\begin{theorem}[Arnold comparison and scalar companion values]
\label{thm:climax-vol1-platonic}
\ClaimStatusConditional
The following statements have the indicated type signatures.
\begin{enumerate}[label=\textup{(\roman*)}]
\item \emph{Arnold window.}  For every datum
$(A;V,n,H_{\mathrm{KZ}})$, the representation
$\widetilde\rho_{A,V,n}$ defines a functor
\begin{equation}
\label{eq:KZ-functor-signature-platonic}
 \KZ_{(-;V,n)}\colon
 \ChirAlg^{H_{\mathrm{KZ}}}_{\mathbb P^1;V,n}
 \longrightarrow
 \ConnConf^{\mathrm{Arn,log}}_{\mathbb P^1,n}
\end{equation}
and Equation~\eqref{eq:climax-equation-G2}.  The
Fulton--MacPherson residue realization is the bar-differential
comparison.

\item \emph{Scalar companion values.}  On the named companion pairs
and hypothesis packages of Chapter~\ref{ch:landscape-census},
\begin{equation}
\label{eq:climax-scalar-companion-values}
 K^\kappa(A)\in\left\{0,13,\frac{250}{3}\right\}.
\end{equation}
The value $0$ occurs on the Heisenberg, affine, and
$\beta\gamma/bc$ pairs; $13$ occurs on the Virasoro pair; and
$250/3$ occurs on the principal $\mathcal W_3$ pair.

\item \emph{Bershadsky--Polyakov central identity.}  In the
Fehily--Kawasetsu--Ridout conformal convention,
\begin{equation}
\label{eq:climax-bp-central-lane}
c_{\mathrm{BP}}(k)
=-\frac{(2k+3)(3k+1)}{k+3},
\qquad
K^c_{\mathrm{BP}}=c_{\mathrm{BP}}(k)+c_{\mathrm{BP}}(-k-6)=50.
\end{equation}
The generators $J,G^+,G^-,T$ are even.  The value of
$K^\kappa_{\mathrm{BP}}$ has status \ClaimStatusOpen{} pending the
genus-one curvature calculation of
Proposition~\ref{prop:bp-kappa}.

\item \emph{Mukai lattice invariant.}  The integral Mukai lattice has
signature $(4,20)$, hence
\begin{equation}
\label{eq:climax-mukai-lattice-lane}
K^{\mathrm{lat}}\bigl(\widetilde H(K3,\mathbb Z)\bigr)
=2c_+=8.
\end{equation}
The chiral bridge
$K^\kappa(A_{\mathrm{Muk}})=8$ has status
\ClaimStatusConjectured{} under
$H_{\mathsf B}=(H_{\mathrm{chart}},H_{\mathrm{KD}},H_{\mathrm{scalar}},
H_{\mathrm{mod}},H_{\mathrm{quantum}})$.

\item \emph{Ghost contribution.}  Every equality involving
$c_{\mathrm{gh}}(R_\bullet)$ is indexed by the BRST presentation and
its normalization $\alpha_R$.  Its exact theorem statement is
$\Gamma_{\mathrm{gh}}(A,R_\bullet;\alpha_R)=0$ on the family where
that comparison has been computed.

\item \emph{Five theorem pointers.}  Theorems A, B, C, D, H retain
the separate ambients and hypothesis packages recorded in
Theorems~\ref{thm:theorem-a-platonic-pointer}--%
\ref{thm:theorem-h-platonic-pointer}.  The present theorem supplies
notation and comparison data for those statements.
\end{enumerate}
\end{theorem}

\begin{proof}
\label{prf:climax-vol1-platonic}
Clause~(i) is the enveloping-algebra universal property together with
$H_{\mathrm{Borch}}$ and $H_{\mathrm{FM}}$.  Clause~(ii) is the
family-by-family computation of
Chapter~\ref{ch:landscape-census}.  Clause~(iii) is
Theorem~\ref{thm:bp-koszul-conductor-polynomial} and
Proposition~\ref{prop:bp-kappa}.  Clause~(iv) is the integral-lattice
calculation of
Proposition~\ref{prop:lattice-k3-mukai-lane-separation}; the five
unproved comparison steps are exactly the components of
$H_{\mathsf B}$ in
Proposition~\ref{prop:G-B-heisenberg-rho-bifurcation}.  Clause~(v) is
the definition of $\Gamma_{\mathrm{gh}}$.  Clause~(vi) is the pointer
block below.
\end{proof}

\section{Verification architecture for the Arnold window}
\label{sec:climax-platonic-proof-architecture}

Four classical constructions illuminate different parts of the
Arnold window.  Fulton--MacPherson residues provide the chain map;
Kohno identifies the affine monodromy; Drinfeld supplies associator
coherence; the direct commutator computation verifies flatness after
the infinitesimal braid representation has been constructed.  They
overlap on a common finite window, and each proves its own link in
the comparison chain.

\subsection{Path~A: Fulton--MacPherson residues}
\label{ssec:climax-platonic-path-A-fm}

The first path is purely geometric. The chiral bar differential is
extracted from the Fulton--MacPherson compactification
$\overline{\mathrm{FM}}_n^{\mathrm{ord}}(\Conf)$ via residue maps at
the codimension-one boundary divisors corresponding to two-point
collisions.

\begin{lemma}[FM residue presentation of the collision differential]
\label{lem:fm-residue-dbar-platonic}
\ClaimStatusConditional
In the logarithmic Fulton--MacPherson model used in
Chapter~\ref{chap:bar-construction}, the collision component of the
ordered chiral bar differential at length $n$ is the signed sum over
pairs $(i,j)$ of the residue map
$\Res_{\Delta_{ij}} \colon \Omega^{n}\bigl(\Confnord(\Conf), A^{\boxtimes n}\bigr)
\to \Omega^{n-1}\bigl(\Conf_{n-1}^{\mathrm{ord}}(\Conf), A^{\boxtimes(n-1)}\bigr)$
applied to the logarithmic OPE residue at $\Delta_{ij}$.  The
augmentation and internal-differential components retain their usual
bar-complex definitions.
\end{lemma}

\begin{proof}[Attribution]
The factorization homology differential is assembled from collision
maps.  The logarithmic compactification identifies the binary
collision map with Poincar\'e residue; see
Chapter~\ref{chap:bar-construction} for the model and
Beilinson--Drinfeld (2004, \S4.2) for the factorization construction.
\end{proof}

\begin{lemma}[Pullback presentation under $H_{\mathrm{KZ}}$]
\label{lem:pullback-presentation-nabla-platonic}
\ClaimStatusConditional
Assume $H_{\mathrm{KZ}}(A;V,n)$.  The connection
$\nabla^{B,0}_{A,V,n}$ is the base change of $\nabla_{\Arn}$ along
$\widetilde\rho_{A,V,n}$, and its FM-boundary residue realization is
the collision differential on~$V$.
\end{lemma}

\begin{proof}
$H_{\mathrm{Borch}}$ makes
$t_{ij}\mapsto T_{ij}$ a representation of $\mathfrak t_n$.
Remark~\ref{rem:connconf-initial-platonic} gives the base-change
identity.  $H_{\mathrm{FM}}$ identifies the residue of
$T_{ij}\eta_{ij}$ with the $(i,j)$ collision summand.  Summing over
pairs gives the asserted collision differential.
\end{proof}

\begin{lemma}[Triple-collision commutator; \ClaimStatusConditional]
\label{lem:triple-collision-cybe-platonic}
Under $H_{\mathrm{Borch}}$, the logarithmic collision operators satisfy
\begin{equation}
\label{eq:cybe-on-r}
 [T_{ij},T_{ik}+T_{jk}]=0,
\end{equation}
for every triple of distinct indices.
\end{lemma}

\begin{proof}
This is the shared-index clause of $H_{\mathrm{Borch}}$.  In the
flatness calculation, Arnold's three-term relation cancels the
logarithmic two-form coefficients, while the simple-pole projection of
the Borcherds identity cancels the operator coefficients.
\end{proof}

\begin{remark}[Closure of Path~A]
\label{rem:path-A-closure-platonic}
Lemmas~\ref{lem:fm-residue-dbar-platonic} and
\ref{lem:pullback-presentation-nabla-platonic} establish
Equation~\eqref{eq:climax-equation-G2} for the finite window carrying
$H_{\mathrm{KZ}}$.  The residue-realization functor then recovers its
collision differential.
\end{remark}

\subsection{Path~B: Kohno monodromy}
\label{ssec:climax-platonic-path-B-kohno}

The second path is representation-theoretic and computes the
monodromy of $\nabla(A)$ in evaluation modules.

\begin{theorem}[Kohno 1987]
\label{thm:kohno-monodromy-platonic}
\ClaimStatusProvedElsewhere
For $A = V^{k}(\fg)$ with $\fg$ simple and $k\neq-h^\vee$, the
monodromy of the KZ connection on evaluation modules
$V_1 \otimes \cdots \otimes V_n$ is the quantum-group braid
representation of $U_{q}(\fg)$ at
$q=\exp(\pi i/(k+h^\vee))$ in the convention $K_i=q^{h_i}$.
The convention $q_{\mathrm{alt}}=q^2$ gives the frequently used
$\exp(2\pi i/(k+h^\vee))$ parameter.
\end{theorem}

\begin{proof}[Attribution]
Kohno (\textit{Annales Inst.\ Fourier}~$37$, $1987$) for the genus-zero
KZ monodromy; Drinfel'd (\textit{Algebra i Analiz}~$2$, $1989$) for
the associator construction; Kazhdan--Lusztig
(\textit{J.\ Amer.\ Math.\ Soc.}~$6$--$7$, $1993$--$1994$) for the
type-A category equivalence. See also Vol~I Section~\ref{sec:derived-dk}.
\end{proof}

\begin{remark}[Closure of Path~B for affine input]
\label{rem:path-B-closure-platonic}
Theorem~\ref{thm:kohno-monodromy-platonic} computes the monodromy of
the affine Arnold window.  It verifies the representation-theoretic
output of Equation~\eqref{eq:climax-equation-G2}; the FM residue
comparison remains the chain-level input supplied by
$H_{\mathrm{FM}}$.
\end{remark}

\subsection{Path~C: Drinfeld associator}
\label{ssec:climax-platonic-path-C-drinfeld}

The third path is analytic and computes Arnold's monodromy via
iterated integrals. The Drinfeld associator $\Phi_{\KZ} \in U(\mathfrak{t}_3)
[[\hbar]]$ is the universal monodromy element for $\nabla_{\Arn}$ at
$n = 3$, controlling the failure of strict associativity at triple
collisions.

\begin{theorem}[Drinfeld $1989$]
\label{thm:drinfeld-associator-platonic}
\ClaimStatusProvedElsewhere
The monodromy of $\nabla_{\Arn}$ on $\Conf_3^{\mathrm{ord}}(\mathbb{C})$
through the iterated-integral construction equals the Drinfeld
associator $\Phi_{\KZ}$, satisfying the pentagon and hexagon coherence
identities and the duality relation
$\Phi_{\KZ}(t_{12}, t_{23})\,\Phi_{\KZ}(t_{23}, t_{12}) = 1$.
\end{theorem}

\begin{proof}[Attribution]
Drinfel'd (\textit{Algebra i Analiz}~$2$, $1989$, $\S 4$). See also
Bar-Natan (\textit{Selecta Math.}~$4$, $1998$) for an explicit
combinatorial presentation.
\end{proof}

\begin{remark}[Closure of Path~C]
\label{rem:path-C-closure-platonic}
For every $\mathfrak t_3$ representation, base change sends
$\Phi_{\KZ}$ to the associativity constraint of its KZ monodromy.
Thus the associator controls the triple-collision coherence of every
finite window satisfying $H_{\mathrm{KZ}}$.
\end{remark}

\subsection{Path~D: direct chain-level Yang--Baxter}
\label{ssec:climax-platonic-path-D-direct}

The fourth calculation expands the square of the logarithmic
superconnection and isolates its codimension-two components.

\begin{proposition}[Codimension-two flatness criterion]
\label{prop:cybe-equivalence-platonic}
\ClaimStatusConditional
Let $(A;V,n)$ satisfy $H_{\mathrm{fin}}$ and $H_{\log}$.  The
connection~\eqref{eq:nabla-of-A} is flat precisely when the operators
$T_{ij}$ satisfy the infinitesimal braid relations.  Under
$H_{\mathrm{FM}}$, this flatness calculation maps to the
codimension-two collision component of $\dbar^2|_V=0$.
\end{proposition}

\begin{proof}
Expand the square of
$d_V+d_{\mathrm{dR}}-\sum T_{ij}\eta_{ij}$.  The dg-intertwining
condition removes the terms containing $d_V$.  Disjoint divisors give
$[T_{ij},T_{kl}]$, killed by the disjoint-pair infinitesimal
braid relation; shared-index divisors give
$[T_{ij},T_{ik}+T_{jk}]$ after Arnold's three-term form identity.
The simple-pole projection of the Borcherds identity supplies the
operator relation at a shared-index collision.  These coefficients
vanish exactly under the infinitesimal braid relations.
$H_{\mathrm{FM}}$ transports the calculation to the collision
summand of the bar square.
\end{proof}

\section{Classical comparison theorems}
\label{sec:climax-platonic-classical-specialisations}

The Arnold window directly contains Drinfeld--Kohno.  Verlinde and
Borcherds belong to additional modular and automorphic structures;
their precise relation to the ordered bar is recorded separately.

\subsection{Drinfeld--Kohno: affine Kac--Moody at evaluation modules}
\label{ssec:climax-platonic-cor-DK}

\begin{corollary}[Affine Arnold window and Drinfeld--Kohno]
\label{cor:climax-drinfeld-kohno-platonic}
\ClaimStatusProvedElsewhere
For $A = V^{k}(\fg)$ at generic level $k$ and $\fg$ simple, the
monodromy representation of $\nabla_{\mathrm{KZ}}$ on evaluation modules
$V_1, \dots, V_n$ equals the quantum-group braid representation
$\rho_n^{U_q(\fg)}$ at $q=\exp(\pi i/(k+h^\vee))$ in the
$K_i=q^{h_i}$ convention:
\begin{equation}
\label{eq:climax-drinfeld-kohno-platonic}
 \rho_n^{\KZ}(V_1,\dots,V_n)
 \;=\; \rho_n^{U_q(\fg)}(V_1,\dots,V_n).
\end{equation}
\end{corollary}

\begin{proof}
By Computation~\ref{ex:kz-of-vk-g-platonic}, the connection coefficient
in KZ normalisation is $\Omega_{ij}/((k+\hv)\,z_{ij})$;
hence the displayed connection is the standard KZ connection. By
Theorem~\ref{thm:kohno-monodromy-platonic}, its monodromy on
evaluation modules is $\rho_n^{U_q(\fg)}$ at the displayed~$q$.
\end{proof}

\subsection{Verlinde: rational chiral algebra at genus zero}
\label{ssec:climax-platonic-cor-Verlinde}

\begin{corollary}[Verlinde formula on a modular tensor category]
\label{cor:climax-verlinde-platonic}
\ClaimStatusProvedElsewhere
Let $A$ satisfy the rationality, $C_2$-cofiniteness, self-contragredience,
and semisimplicity hypotheses that make its module category modular.
Then the
fusion coefficients $N_{ij}^{k}$ are diagonalised by $S$:
\begin{equation}
\label{eq:climax-verlinde-platonic}
 N_{ij}^{k}
 \;=\; \sum_{\ell}
 \frac{S_{i\ell}\,S_{j\ell}\,(S^{-1})_{\ell k}}{S_{0\ell}}.
\end{equation}
For a unitary modular tensor category,
$(S^{-1})_{\ell k}=\overline{S_{k\ell}}$.
\end{corollary}

\begin{proof}
This is the Verlinde theorem for the modular tensor category of
$A$-modules.  The genus-zero braid connection and the genus-one
$S$-matrix are joined by the sewing and modularity theorems of the
rational theory; those hypotheses supply the additional structure
beyond $H_{\mathrm{KZ}}$.
\end{proof}

\subsection{Borcherds: lattice vertex algebra at genus one}
\label{ssec:climax-platonic-cor-Borcherds}

\begin{remark}[Borcherds comparison obligation]
\label{cor:climax-borcherds-platonic}
\ClaimStatusOpen
Borcherds' multiplicative lift constructs automorphic products from
vector-valued modular or Jacobi input.  A comparison with
$H^\bullet(\BarOrd(V_\Lambda))$ requires a trace theorem identifying
the relevant Fourier coefficients with the Euler characteristics of
the completed bar complex.  That trace theorem is the named open
obligation for this comparison.  The proved automorphic statement remains
the Borcherds product theorem itself.
\end{remark}

\subsection{Arnold's KZ-monodromy theorem as the common root}
\label{ssec:climax-platonic-cor-Arnold}

\begin{corollary}[Arnold universal coefficient map]
\label{cor:climax-arnold-common-root-platonic}
\ClaimStatusProvedHere
Every finite window satisfying $H_{\mathrm{KZ}}$ carries the algebra
map
\begin{equation}
\label{eq:four-shadows-platonic}
 U(\mathfrak t_n)\longrightarrow\End(V),
 \qquad t_{ij}\longmapsto T_{ij},
\end{equation}
and its monodromy is the corresponding specialization of Arnold's
pure-braid representation.
\end{corollary}

\begin{proof}
This is Remark~\ref{rem:connconf-initial-platonic} applied to
$\widetilde\rho_{A,V,n}$.
\end{proof}

\section{Five Vol~I theorem pointers}
\label{sec:climax-platonic-five-theorems-corollaries}

Theorems A, B, C, D, H retain distinct inputs and reconstruction
maps.  The present chapter fixes the common notation used by their
comparison maps.  A concerns the bar--cobar adjunction; B concerns
strict and completed inversion; C concerns Verdier complementarity and
its typed scalar projections; D concerns modular pushforward; H
concerns the chiral derived centre.  For a finite set
$S\subset\mathbb Z$, the Theorem~H package $H_H(\cA;S)$ consists of a
complete chart-cochain model and incidence-compatible strong
deformation retract onto a complex supported in $S$.  The pointers
below name the chapters where each theorem is stated and proved in
its native form; the corresponding pointer remark identifies the
categorical direction being read.

\begin{theorem}[Theorem~A: chiral bar--cobar adjunction, pointer]
\label{thm:theorem-a-platonic-pointer}
\ClaimStatusConditional
\textup{[Type signature: Open quadrant; augmented factorization
presentation; Beilinson levels \(1\leftrightarrow2\);
pro-nilpotent chiral Ran ambient;
\(H_A=H_{\mathrm{fact}}\cup H_{\mathrm{conv}}\).]}
On the stable factorization ambient specified by \(H_A\), chiral
cobar and bar form the adjunction
\[
 \boldsymbol\Omega_X
 :
 \mathsf{FactCoalg}_X^{\mathrm{conil,comp}}
 \rightleftarrows
 \mathsf{FactAlg}_X^{\mathrm{aug}}
 :
 \mathbf B_X,
 \qquad
 \boldsymbol\Omega_X\dashv\mathbf B_X.
\]
Its chain realization has counit
\[
 \psi_\cA\colon
 \Omega^{\mathrm{ch}}\bar B^{\mathrm{ch}}(\cA)\longrightarrow\cA
\]
induced by the universal twisting morphism.  On the quadratic chiral
presentation, let
$\tau_{\mathrm i}\colon\cA^i\to\cA$ be the canonical twisting
morphism.  Priddy's coalgebra carries the additional map
\[
 q_\cA\colon\cA^i\longrightarrow B_X(\cA),
\]
which is an equivalence under
$H_{\mathrm{CL}}(\cA,\cA^i,\tau_{\mathrm i})$.  This quadratic
comparison is separate from the full bar--cobar unit and counit.
\end{theorem}

\begin{proposition}[Chain and presentable-\((\infty,1)\) readouts]
\label{prop:theorem-a-dual-lane-local-readout}
\ClaimStatusConditional
Under \(H_A\), the chain counit \(\psi_\cA\) is the model-level
realization of the presentable-\((\infty,1)\) counit
\(\boldsymbol\Omega_X\mathbf B_X(\cA)\to\cA\).
The model/realization package supplies a natural comparison between
these maps and asserts its compatibility with the Ran completion and
factorization structure.
\end{proposition}

\begin{proof}
The universal twisting morphism gives the chain adjunction.  The
enhanced bar--cobar theorem in the pro-nilpotent Ran ambient gives the
categorical counit, and the clauses of $H_{\mathrm{fact}}$ restrict it
to the factorization subcategory.  Naturality of the realization
comparison identifies the two representatives.
\end{proof}

\begin{theorem}[Theorem~B: strict and completed inversion, pointer]
\label{thm:theorem-b-platonic-pointer}
\ClaimStatusConditional
\textup{[Type signature: Open quadrant; chiral CDG presentation;
Beilinson level \(2\); strict package \(H_B^{\mathrm{str}}\) or
completed package \(H_B^{\mathrm{comp}}\).]}
Theorem~B has two theorem surfaces.
\begin{enumerate}[label=\textup{(\alph*)}]
\item Under \(H_B^{\mathrm{str}}\), the counit
\[
 \Omega_X^{\mathrm{ch}}\bar B_X^{\mathrm{ch}}(\cA)
 \longrightarrow \cA
\]
is a chain quasi-isomorphism
(Theorems~\ref{thm:higher-genus-inversion} and
\ref{thm:bar-cobar-inversion-qi}).

\item Under \(H_B^{\mathrm{comp}}\), let
\(\widehat C=\widehat{\bar B}^{\mathrm{ch}}(\cA)\) be the
weight-complete conilpotent chiral CDG coalgebra.  The continuous
co-/contra-derived comparison is
\[
 D^{\mathrm{co}}\!\left(
 \widehat C\text{-}\mathrm{comod}^{\mathrm{conil,ch}}\right)
 \simeq
 D^{\mathrm{ctr}}\!\left(
 \widehat C\text{-}\mathrm{contra}^{\mathrm{ch}}\right)
\]
(Theorem~\ref{thm:chiral-positselski-weight-completed}).
The curved structure is interpreted in the second-kind derived
categories, and its comparison morphisms are inverted through
coacyclic cones.
\end{enumerate}
The strict and completed packages specify distinct categories and
distinct notions of equivalence.
\end{theorem}


\begin{theorem}[Theorem~C: complementarity and typed scalar projections]
\label{thm:theorem-c-platonic-pointer}
\ClaimStatusConditional
Theorem~C has a strict-flat centre-local-system core and a
family-indexed scalar readout.
\begin{enumerate}[label=\textup{(\alph*)}]
\item \emph{Centre-local-system complementarity.}
Under the C0/C1 hypothesis package of
Theorem~\ref{thm:quantum-complementarity-main}, Verdier duality gives
the eigenspace decomposition and Lagrangian pairing
\[
 \mathbf C_g(\cA)
 \simeq
 \mathbf Q_g(\cA)\oplus\mathbf Q_g(\cA^!).
\]
Here
$\mathbf C_g(\cA)=R\Gamma(\overline{\mathcal M}_g,\mathcal Z(\cA))$
uses the ordinary centre local system obtained from the strict-flat C0
comparison.  The Theorem~H derived centre connects to this coefficient
system through a supplied $\iota_Z^{\mathrm{der}}$, followed by~$H^0$
and C0.
The shifted-symplectic BV extension carries the additional C2 package
of Theorem~\ref{thm:shifted-symplectic-complementarity}.

\item \emph{Computed modular conductors.}
For the named companion pairs in the canonical census,
\[
 K^\kappa(\cA)=\kappa(\cA)+\kappa(\cA^!)
 \in
 \left\{0,\ 13,\ \frac{250}{3}\right\}.
\]
The three values are realized by the Heisenberg/affine/free-field
pairs, the Virasoro pair, and the principal \(\mathcal W_3\) pair,
respectively.  Each row carries the duality package displayed in
Chapter~\ref{ch:landscape-census}.

\item \emph{Bershadsky--Polyakov.}
The standard Fehily--Kawasetsu--Ridout conformal vector gives the
proved central-charge conductor
\[
 K^c_{\mathrm{BP}}
 =c_{\mathrm{BP}}(k)+c_{\mathrm{BP}}(-k-6)=50.
\]
The generators \(J,G^+,G^-,T\) are even.  The modular entries
\(\kappa_{\mathrm{BP}}\) and \(K^\kappa_{\mathrm{BP}}\) have status
\ClaimStatusOpen{} pending the direct genus-one curvature computation
of Proposition~\ref{prop:bp-kappa}.  Secondary conformal-vector
normalizations form a separate family whose genus-one transformation
law remains to be proved.

\item \emph{Mukai lattice candidate.}
The proved lattice invariant is
\[
 K^{\mathrm{lat}}\bigl(\widetilde H(K3,\mathbb Z)\bigr)
 =2c_+=8.
\]
Its interpretation as
\(K^\kappa(A_{\mathrm{Muk}})=8\) has status
\ClaimStatusConjectured{} under
\[
 H_{\mathsf B}
 =(H_{\mathrm{chart}},H_{\mathrm{KD}},H_{\mathrm{scalar}},
   H_{\mathrm{mod}},H_{\mathrm{quantum}}).
\]
Equation~\eqref{eq:climax-scalar-companion-values} therefore contains
three computed values; adjoining the \(H_{\mathsf B}\)-candidate gives
the conditional numerical set \(\{0,8,13,250/3\}\).
\end{enumerate}
\end{theorem}

\begin{remark}[The Mukai arithmetic and its chiral bridge]
\label{rem:climax-k3-face}
The integral Mukai lattice of a complex K3 surface is
\[
 \widetilde H(K3,\mathbb Z)
 =
 H^0(K3,\mathbb Z)\oplus H^2(K3,\mathbb Z)\oplus H^4(K3,\mathbb Z)
 \simeq U^4\oplus E_8(-1)^2.
\]
It has rank \(24\) and signature \((4,20)\).  Consequently
\(c_+=4\) and \(2c_+=8\).  This calculation is a theorem of integral
lattice theory.

The proposed chiral reading requires all five maps in
\(H_{\mathsf B}\).  Here \(H_{\mathrm{chart}}\) constructs an augmented
chiral chart from the indefinite lattice,
\(H_{\mathrm{KD}}\) proves its bar comparison and identifies the
Koszul companion, \(H_{\mathrm{scalar}}\) identifies the modular pair
sum with \(2c_+\), \(H_{\mathrm{mod}}\) constructs the
Humbert--Borcherds comparison, and \(H_{\mathrm{quantum}}\) constructs
a functor selecting a root order.  Thus the equation
\(K^\kappa(A_{\mathrm{Muk}})=8\) is a precise five-obligation
conjecture.

Bruinier (2002, Lemma~5.1) supplies a coefficient-field statement for
vector-valued modular forms, with \(\operatorname{lcm}(N,8)\) in the
half-integral-weight case.  Lusztig (1990) constructs quantum groups
at a chosen admissible root order.  These are ambient inputs for
\(H_{\mathrm{mod}}\) and \(H_{\mathrm{quantum}}\); the comparison maps
themselves remain the open mathematics.
\end{remark}

\begin{remark}[K3 Hall--Drinfeld recognition obligations]
\label{rem:k3-koszul-structure-precise}
\ClaimStatusOpen
Let \(\mathbf H_{\Delta_5}\) denote the proposed K3
Hall--Drinfeld recognition target of Volume~III.  A theorem placing it
on the strict chiral Koszul surface requires the following data in one
ambient:
\begin{enumerate}[label=\textup{(\alph*)}]
\item a constructed factorization or chiral algebra with specified
generators, relations, parity, completion, and coproduct;
\item a PBW theorem determining the full relation ideal of the
Hall/CoHA presentation;
\item a conilpotent or weight-complete bar model and an identified
Verdier companion;
\item an explicit twisting morphism whose bar--cobar counit is a
quasi-isomorphism on a named locus;
\item a modular comparison relating the chiral family to
\(\Delta_5\), with nearby-cycle data at every asserted Humbert wall.
\end{enumerate}
These five constructions form the recognition package
\(H_{\Delta_5}^{\mathrm{rec}}\).  The BKM denominator theorem for
\(\Delta_5\) and the Mukai lattice calculation supply arithmetic
targets; \(H_{\Delta_5}^{\mathrm{rec}}\) supplies the algebraic bridge.
The present status is \ClaimStatusOpen.
\end{remark}


\begin{theorem}[Theorem~D: deformation, Hodge, and graph traces; pointer]
\label{thm:theorem-d-platonic-pointer}
\ClaimStatusConditional
\textup{[Type signature: Open quadrant; modular-trace and cyclic-graph
presentations; Beilinson levels \(4\to5\); hypothesis package
\(H_D=(H_D^1,H_D^K,H_D^{\mathrm{tr}},H_D^{\mathrm{graph}})\).]}
For \(g\geq2\), the native deformation class is
\[
 \operatorname{Obs}^{\mathrm{def}}_g(\cA)
 \in H^2\!\left(\Def_g(\cA)\right).
\]
Under~\(H_D^1(\cA)\), the pointed genus-one class has trace
\[
 \operatorname{tr}_1\operatorname{Obs}^{\mathrm{def}}_{1,1}(\cA)
 =\kappa(\cA)\lambda_1.
\]
Under~\(H_D^K(\cA,g)\), the scalar virtual object and its Hodge
image are
\[
 \mathfrak O_g^K(\cA)=\kappa(\cA)\lambda_{-1}(\mathbb E_g),
 \qquad
 \operatorname{ch}_g\!\left(\mathfrak O_g^K(\cA)\right)
 =(-1)^g\kappa(\cA)\lambda_g.
\]
The package~\(H_D^{\mathrm{tr}}\) supplies the derived comparison
from the degree-two deformation class to this perfect virtual object.
Under~\(H_D^{\mathrm{graph}}(\cA,g)\), the all-weight numerical trace of
Theorems~\ref{thm:modular-characteristic} and
\ref{thm:multi-weight-genus-expansion} satisfies
\[
 F_g(\cA)
 =
 \kappa(\cA)\lambda_g^{\mathrm{FP}}
 +
 \delta F_g^{\mathrm{cross}}(\cA).
\]
Here \(\delta F_1^{\mathrm{cross}}=0\); the scalar-diagonal graph
package gives \(\delta F_g^{\mathrm{cross}}=0\) in every genus.
\end{theorem}

\begin{theorem}[Theorem~H: chiral derived centre, pointer]
\label{thm:theorem-h-platonic-pointer}
\ClaimStatusConditional
\textup{[Type signature: Open/Bulk comparison; chiral Hochschild
presentation; Beilinson level \(3\); hypothesis package
\(H_H(\cA;S)\).]}
The chart cochain object is
\[
 \ChirHoch^\bullet(\cA)
 =
 C^\bullet_{\mathrm{ch}}(\cA,\cA)
 =
 Z^{\mathrm{der}}_{\mathrm{ch}}(\cA)
\]
in the chart presentation.  Fix a finite set $S\subset\mathbb Z$.
The package $H_H(\cA;S)$ constructs a complete filtered model and an
incidence-compatible strong deformation retract onto an
$S$-supported complex.  It gives
\[
 \operatorname{Supp}\ChirHoch^\bullet(\cA)\subseteq S.
\]
Bakalov--De Sole--Kac compute the bounded benchmarks
$\{0,2,3\}$ for Virasoro and dimension vector $(2,1,0,\ldots)$ for
the rank-one even free superboson.  Transport from bounded vertex
cochains to the chiral chart uses a separately named quasi-isomorphism.
Topological and categorical Hochschild theories retain their typed
comparison maps.
\end{theorem}

\begin{remark}[Five theorem packages in one coordinate chart]
\label{rem:five-corollaries-together-platonic}
Theorems~\ref{thm:theorem-a-platonic-pointer}--%
\ref{thm:theorem-h-platonic-pointer} use the same Beilinson levels and
companion notation.  Their reconstruction maps are, respectively, the
bar--cobar unit/counit, strict or completed inversion, Verdier
complementarity, modular pushforward, and the chiral derived-centre
comparison.  The common coordinate chart preserves the five native
proofs.
\end{remark}


\section{Open comparison obligations \(\Pi_1\)--\(\Pi_4\)}
\label{sec:climax-platonic-open-obstructions}

The four obligations below enlarge a proved input by one named
construction.  Each is stated in the ambient where that construction
must live.

\begin{conjecture}[\(\Pi_1\): elliptic and higher-genus residue comparison]
\label{conj:climax-kzb-genus-one-platonic}
\ClaimStatusConjectured
For an elliptic curve \(E_\tau\), let
\begin{equation}
\label{eq:nabla-kzb-platonic}
 \nabla_{\mathrm{KZB}}
 =
 d-\sum_{i<j}t_{ij}\,\eta^{(1)}_{ij}(\tau,z_{ij})
 +\text{the standard dynamical/Cartan term}
\end{equation}
denote an elliptic KZB connection in a fixed convention.  An elliptic
analogue of Equation~\eqref{eq:climax-equation-G2} requires a package
\(H_{\mathrm{KZB}}\) consisting of an elliptic bar model, a global
propagator and dynamical term, collision operators satisfying the
elliptic braid relations, and a residue comparison compatible with
degeneration.  Calaque--Enriquez--Etingof construct universal KZB
connections; \(H_{\mathrm{KZB}}\) is the additional chiral-bar
comparison.
\end{conjecture}

\begin{conjecture}[\(\Pi_2\): principal-\(\mathcal W\) collision operators]
\label{conj:climax-w-arena-platonic}
\ClaimStatusConjectured
For a principal \(\mathcal W^k(\mathfrak g)\), a KZ-type bar window
requires explicit logarithmic collision operators on a finite bar
model and a proof of the infinitesimal braid relations from the
screening or BRST presentation.  The package
\(H_{\mathcal W\mathrm{KZ}}\) consists of those operators, their
functoriality, and the FM residue comparison.  Principal screening
complexes provide the natural source of the construction.
\end{conjecture}

\begin{conjecture}[\(\Pi_3\): functorial BRST comparison]
\label{conj:climax-universal-brst-platonic}
\ClaimStatusConjectured
Let \(\mathcal C_{\mathrm{BRST}}\) be a specified subcategory of chiral
algebras equipped with chosen reduction data.  The desired
construction is a functor
\[
 \mathsf R_{\mathrm{BRST}}\colon
 \mathcal C_{\mathrm{BRST}}\longrightarrow
 \mathsf{BRSTGaugedChirAlg}
\]
together with a natural normalization
\(\alpha_R\) and a proof that
\(\Gamma_{\mathrm{gh}}(A,R_\bullet;\alpha_R)=0\) on each connected
family.  Non-principal reductions require the charged, neutral, and
conformal-improvement sectors in the same genus-one calculation.
Bershadsky--Polyakov is the first open test.
\end{conjecture}

\begin{conjecture}[\(\Pi_4\): chiral--Borcherds trace comparison]
\label{conj:climax-universal-trace-identity-platonic}
\ClaimStatusConjectured
For a Borcherds product \(\Phi_N\), the automorphic theorem gives
\[
 \kappa_{\mathrm{BKM}}(\Phi_N)
 =\operatorname{wt}(\Phi_N)
 =\frac{c_N(0)}2.
\]
For a CY-to-chiral output \(A_X\), the proposed comparison asks for an
isomorphism of determinant lines
\begin{equation}
\label{eq:trace-identity-platonic}
 \Theta_{\Phi}\colon
 \mathcal L_{\mathrm{ch}}(A_X)
 \xrightarrow{\ \sim\ }
 \mathcal L_{\mathrm{Bor}}(\Phi_N),
\end{equation}
whose equality of weights compares \(K^\kappa(A_X)\) with
\(\kappa_{\mathrm{BKM}}(\Phi_N)\).  Its hypothesis package
\(H_{\mathrm{tr}}\) consists of the CY-to-chiral construction,
chiral Koszulness of \(A_X\), the determinant-line comparison, and
compatibility with the Borcherds lift.  The Mukai lattice number \(8\),
the Borcherds weight \(5\) of \(\Delta_5\), and every BRST ghost charge
remain separate entries of~\eqref{eq:climax-equation-G3} until
\(H_{\mathrm{tr}}\) is proved.
\end{conjecture}

\section{Verification anchors}
\label{sec:climax-platonic-iv}

\subsection{Arnold window}
\label{ssec:climax-platonic-iv-G2}

The verification chain has four typed links.
\begin{enumerate}
\item Arnold's three-term form relation and the infinitesimal braid
relations prove flatness of the universal logarithmic connection.
\item The Borcherds identity proves the operator relations
\(H_{\mathrm{Borch}}\) for a chosen collision model.
\item The FM boundary map supplies \(H_{\mathrm{FM}}\), relating the
superconnection residue to the bar collision differential.
\item Kohno and Drinfeld compute the monodromy and associator of the
affine instance.
\end{enumerate}
The first two links are the form and operator halves of flatness; the
third is the chain comparison; the fourth is a representation-theoretic
check on a named family.

\subsection{Scalar values and statuses}
\label{ssec:climax-platonic-iv-G3}

The canonical census gives the following status table.
\begin{center}
\begin{tabular}{l|c|c|l}
family & scalar type & value & status \\ \hline
Heisenberg / affine / \(\beta\gamma\)-\(bc\)
 & \(K^\kappa\) & \(0\) & computed on named companion pairs \\
Virasoro & \(K^\kappa\) & \(13\) & computed under \(H_{\mathrm{Vir}}\) \\
principal \(\mathcal W_3\)
 & \(K^\kappa\) & \(250/3\) & computed under \(H_{\mathcal W_3}\) \\
Bershadsky--Polyakov & \(K^c\) & \(50\) & proved central identity \\
Bershadsky--Polyakov & \(K^\kappa\) & \(\mathrm{open}\)
 & genus-one curvature obligation \\
Mukai lattice & \(K^{\mathrm{lat}}\) & \(8\) & proved lattice arithmetic \\
Mukai chiral chart & \(K^\kappa\) & \(8\)
 & conjectural under \(H_{\mathsf B}\) \\
\(\Delta_5\) & \(\kappa_{\mathrm{BKM}}\) & \(5\)
 & proved Borcherds weight
\end{tabular}
\end{center}
The Bershadsky--Polyakov proof uses
Theorem~\ref{thm:bp-koszul-conductor-polynomial} and the even
generating fields of Proposition~\ref{prop:bp-kappa}.  The Mukai proof
uses the decomposition
\(U^4\oplus E_8(-1)^2\) and the signature \((4,20)\).  These
computations are independent because they evaluate different scalar
types.

\section{Typed cross-volume comparisons}
\label{sec:climax-platonic-outlook}

\subsection{Vol~II boundary-linear comparison}
\label{ssec:climax-platonic-volii-bridge}

Under the boundary-linear, tempered, and derived-centre package
\(H_{\mathrm{BL}}\), the Vol~II comparison assigns to \(A\) the pair
\[
 \bigl(Z^{\mathrm{der}}_{\mathrm{ch}}(A),A\bigr)
\]
with its
\(\mathrm{SC}^{\mathrm{ch},\mathrm{top}}\)-action and boundary
restriction
\[
 \rho_\partial\colon Z^{\mathrm{der}}_{\mathrm{ch}}(A)\longrightarrow A.
\]
This map compares the level-three derived centre with the level-one
boundary algebra.  The bar coalgebra remains the level-two object in
the Beilinson tower.

\subsection{Vol~III CY-to-chiral and trace comparison}
\label{ssec:climax-platonic-voliii-bridge}

The proposed Vol~III construction has the two-stage type
\[
 \Phi_d^{(\Sigma_{d-1},C)}
 =
 \mathrm{Sp}^{\mathrm{ch}}_{\Sigma_{d-1},C}
 \circ
 \Phi_d^{\mathrm{FA}}.
\]
Its application in this chapter carries the package \(H_\Phi\):
construction of the factorization algebra
\(\Phi_d^{\mathrm{FA}}(\mathcal C)\), construction of the
specialization functor, and verification of the resulting operadic
level on \(C\).  When \(H_\Phi\) produces a chirally Koszul
\(A_X\), the Vol~I theorem packages apply in their native ambients.
The Borcherds formula
\(\kappa_{\mathrm{BKM}}(\Phi_N)=c_N(0)/2\) is the proved automorphic
target; Conjecture~\ref{conj:climax-universal-trace-identity-platonic}
is the determinant-line bridge.

\subsection{Summary of the four extensions}
\label{ssec:climax-platonic-closing}

The established statements consist of the finite-window Arnold
comparison, the scalar values
\(\{0,13,250/3\}\), the BP central conductor \(50\), and the Mukai
lattice invariant \(8\).  The four extensions add, respectively, an
elliptic residue theorem, principal-\(\mathcal W\) collision
operators, a functorial BRST comparison, and a chiral--Borcherds trace
map.  Each extension has a named hypothesis package and a concrete
first test.

\section{K3 and arithmetic frontier: audited obligations}
\label{sec:cclimax-supplement}

\begin{remark}[K3 Hall--Drinfeld recognition target]
\label{rem:cclimax-1}
\ClaimStatusOpen
The notation \(\mathbf H_{\Delta_5}\) denotes the recognition target
specified by the proposed Hall/CoHA source, the Mukai pairing, and the
Borcherds denominator \(\Delta_5\).  Its construction as a chiral
bialgebra carries the package
\(H_{\Delta_5}^{\mathrm{rec}}\) of
Remark~\ref{rem:k3-koszul-structure-precise}.  The proved inputs are
the Mukai lattice and the automorphic denominator theorem.  PBW,
completion, radical quotient, parity, coproduct, and chiral
factorization are the remaining recognition maps.
\end{remark}

\begin{remark}[\(\hbar\)-normalization at the Mukai candidate]
\label{rem:cclimax-2}
The lattice calculation gives \(K^{\mathrm{lat}}=8\).  Choosing the
formal parameter by \(\hbar^2=-1/8\) yields
\[
 \hbar^2K^{\mathrm{lat}}=-1.
\]
This is a normalization identity.  Its promotion to a chiral quantum
statement requires \(H_{\mathrm{scalar}}\) and
\(H_{\mathrm{quantum}}\): the first identifies
\(K^\kappa\) with \(K^{\mathrm{lat}}\), and the second constructs a
quantization functor whose root order is eight.
\end{remark}

\section{A decoupled Saito--Kurokawa arithmetic check}
\label{sec:cclimax-supplement-ii}

The arithmetic calculation below is retained as an exact local-factor
check.  Its role in this chapter is purely arithmetic.

\begin{proposition}[Saito--Kurokawa spinor coefficients of
\(\Delta_{10}\) for \(p\leq37\)]
\label{prop:cclimax-SK-spinor}
\ClaimStatusProvedHere
Let \(f_{18}=E_6\Delta\) be the normalized weight-\(18\) eigenform.
For the Saito--Kurokawa lift
\(\Delta_{10}=\mathrm{SK}(f_{18})\), the spinor coefficient satisfies
\begin{equation}
\label{eq:cclimax-SK-hecke}
 \lambda_p(\Delta_{10})
 =a_p(f_{18})+p^8+p^9.
\end{equation}
For
\(p=2,3,5,7,11,13,17,19,23,29,31,37\), the pairs
\((a_p,\lambda_p)\) are
\[
\begin{array}{r|r|r}
2&-528&240\\
3&-4\,284&21\,960\\
5&-1\,025\,850&1\,317\,900\\
7&3\,225\,992&49\,344\,400\\
11&-753\,618\,228&1\,818\,688\,344\\
13&2\,541\,064\,526&13\,961\,294\,620\\
17&-5\,429\,742\,318&120\,133\,891\,620\\
19&1\,487\,499\,860&341\,158\,760\,680\\
23&-317\,091\,823\,464&1\,562\,371\,823\,280\\
29&2\,433\,410\,602\,590&17\,440\,802\,991\,420\\
31&-8\,849\,722\,053\,088&18\,442\,791\,145\,024\\
37&12\,691\,652\,946\,662&146\,165\,872\,195\,660
\end{array}
\]
and \(|a_p|\leq2p^{17/2}\).
\end{proposition}

\begin{proof}
The Saito--Kurokawa local spinor factor is
\[
 L_p(s,\Delta_{10},\mathrm{spin})
 =
 \zeta_p(s-8)\zeta_p(s-9)L_p(s,f_{18}).
\]
Its four reciprocal roots are
\(p^8,p^9,\alpha_p,\beta_p\), with
\(\alpha_p+\beta_p=a_p(f_{18})\) and
\(\alpha_p\beta_p=p^{17}\).  Their sum gives
Equation~\eqref{eq:cclimax-SK-hecke}.  Expanding
\(E_6(q)\,q\prod_{n\geq1}(1-q^n)^{24}\) gives the table.  Deligne's
theorem gives the final bound.
\end{proof}

\begin{remark}[Extended coefficient window]
\label{rem:cclimax-ap-extension}
\label{rem:cclimax-ap-extension-ii}
\label{rem:cclimax-ap-extension-iii}
\ClaimStatusProvedHere
The same integral \(q\)-series convolution has been run through the
primes \(p\leq199\).  The proof uses the coefficient
identity for \(E_6\Delta\), the Hecke recursion
\(a_{p^2}=a_p^2-p^{17}\), and Deligne's bound.  The tables formerly
printed in this chapter are computational data; the theorem-level
content is the local-factor formula above.
\end{remark}

\begin{remark}[Local Euler factor]
\label{rem:cclimax-sk-euler}
For every finite prime,
\[
 Z_p(s)=\zeta_p(s-8)\zeta_p(s-9)L_p(s,f_{18})
\]
is regular at \(s=19/2\): its three denominator factors evaluate at
\(p^{-3/2}\), \(p^{-1/2}\), and
\(1-a_pp^{-19/2}+p^{-2}\), and Deligne's bound controls the last.
\end{remark}

\begin{remark}[Normalized Satake identity and the open Casimir bridge]
\label{rem:cclimax-satake-casimir}
\label{rem:cclimax-satake-casimir-ii}
Let
\(\widetilde\alpha_p=\alpha_p/p^{17/2}\) and
\(\widetilde\beta_p=\beta_p/p^{17/2}\).  Then
\[
 \widetilde\alpha_p\widetilde\beta_p=1,\qquad
 \widetilde\alpha_p+\widetilde\beta_p=\frac{a_p}{p^{17/2}},
\]
and the exact degree-two identity is
\begin{equation}
\label{eq:cclimax-casimir-dictionary}
 \widetilde\alpha_p^2+\widetilde\beta_p^2
 =\frac{a_p^2}{p^{17}}-2.
\end{equation}
A Frenkel--Reshetikhin Casimir interpretation on
\(\mathbf H_{\Delta_5}\) requires a constructed representation and a
Satake--chiralization functor.  That comparison has status
\ClaimStatusOpen.
\end{remark}

\begin{remark}[Monster quantization datum]
\label{rem:cclimax-monster-lusztig-conway-norton}
\ClaimStatusOpen
A Lusztig root order is an input to a root-of-unity quantum group.
For a proposed Monster Hall--Drinfeld object, write
\begin{equation}
\label{eq:cclimax-monster-lusztig-level}
 \ell_{\mathrm{Monster}}
 :=\operatorname{ord}(q_{\mathrm{Monster}}).
\end{equation}
A theorem fixing this integer requires a quantization functor from the
Monster vertex-algebraic or Borcherds data to a specified quantum
group.  The Fricke involution of the Hauptmodul has order two; the
quantization functor is a separate input.  The value of
\(\ell_{\mathrm{Monster}}\) therefore has status \ClaimStatusOpen.
\end{remark}

\begin{remark}[Metaplectic component group for \(\Delta_5\)]
\label{rem:cclimax-S-psi-delta5}
\ClaimStatusOpen
For a genuine automorphic representation on a metaplectic cover, the
Arthur-type component group would be defined from a constructed
\(L\)-parameter by
\begin{equation}
\label{eq:cclimax-S-psi-delta5}
 S_{\psi_{\Delta_5}}
 :=
 \pi_0\!\left(
 \operatorname{Cent}(\operatorname{im}\psi_{\Delta_5})/
 Z(\widehat G)^\Gamma
 \right).
\end{equation}
The parameter, dual-group convention, and multiplicity formula form
the required input.  The cardinality of the associated packet
\begin{equation}
\label{eq:cclimax-Psi-delta5-total}
 |\Psi_{\Delta_5}| \quad\text{is an open output of that package.}
\end{equation}
\end{remark}

\section{K3 Koszul and Humbert-wall frontier}
\label{sec:cclimax-supplement-iii}

\begin{remark}[Koszul locus as a moduli problem]
\label{rem:cclimax-koszul-locus}
\ClaimStatusOpen
For a constructed chiral family \(A\) over a base \(S\), define
\begin{equation}
\label{eq:cclimax-koszul-locus-sharp}
 U_{\mathrm{Kosz}}(A)
 :=
 \{s\in S:\Omega^{\mathrm{ch}}B^{\mathrm{ch}}(A_s)\to A_s
 \text{ is a quasi-isomorphism}\}.
\end{equation}
Identifying this open locus with a complement of Humbert divisors
requires a family \(A/S\), a perfect bar complex, and a degeneracy
calculation for its counit.  These data belong to
\(H_{\Delta_5}^{\mathrm{rec}}\) and supplement Bruinier's
coefficient-field lemma.
\end{remark}

\begin{lemma}[The \(H_1\) nearby-cycle obligation]
\label{lem:cclimax-H1-failure}
\ClaimStatusOpen
A theorem at \(H_1\) requires the nearby-cycle complex
\(\psi_{H_1}B^{\mathrm{ch}}(A)\), its monodromy operator, and a proof
that the completed bar filtration satisfies the required
Mittag--Leffler condition.  The order of any finite monodromy is an
output of this calculation.
\end{lemma}

\begin{lemma}[The \(H_3\) obstruction-class obligation]
\label{lem:cclimax-H3-failure}
\ClaimStatusOpen
A theorem at \(H_3\) requires a Maurer--Cartan model for the family, an
explicit class in the deformation complex of the bar--cobar counit,
and a homotopy whose differential equals that class.  A physical
three-loop coefficient becomes part of this statement after a
comparison theorem identifies the physical determinant line with the
chiral deformation line.
\end{lemma}

\begin{lemma}[The \(H_4\) descent obligation]
\label{lem:cclimax-H4-failure}
\ClaimStatusOpen
A theorem at \(H_4\) requires a local system on the constructed chiral
family, its étale descent cocycle, and a comparison between that
cocycle and the completed bar complex.  Kuga--Satake and Torelli
theorems provide geometric inputs after the relevant moduli map has
been constructed.
\end{lemma}

\begin{theorem}[Conditional global inversion template]
\label{thm:climax-W18-global-inversion}
\ClaimStatusConditional
Let \(A/S\) be a chiral family satisfying
\(H_{\Delta_5}^{\mathrm{rec}}\).  Assume:
\begin{enumerate}[label=\textup{(\alph*)}]
\item the bar--cobar counit is a quasi-isomorphism on
\(U_{\mathrm{Kosz}}(A)\);
\item every boundary stratum carries a monodromy-refined complete
filtration with a coacyclic counit cone;
\item the local completed inversions satisfy derived Čech descent.
\end{enumerate}
Then the counit is an isomorphism in the resulting global coderived
category.  Strict chain-level inversion holds on
\(U_{\mathrm{Kosz}}(A)\).
\end{theorem}

\begin{proof}
Clause (a) gives the open part, clause (b) gives the formal
neighbourhoods, and clause (c) glues the two in the coderived category.
The mathematical work for \(\mathbf H_{\Delta_5}\) is the
construction and verification of these three hypotheses.
\end{proof}

\begin{remark}[Physical three-loop comparison]
\label{rem:cclimax-W17-costello-link}
A coefficient in a twisted-supergravity amplitude and a class
obstructing a chiral bar--cobar counit live in distinct deformation
complexes.  The comparison map between those complexes is the
load-bearing theorem required to identify them.
\end{remark}

\begin{remark}[Theorem~B scope on the K3 target]
\label{rem:cclimax-theoremB-scope-update}
Theorem~B applies to \(\mathbf H_{\Delta_5}\) after
\(H_{\Delta_5}^{\mathrm{rec}}\) constructs the object and verifies the
strict or completed hypotheses of
Theorem~\ref{thm:climax-W18-global-inversion}.  The current K3
statement is therefore a conditional application template.
\end{remark}

\section{Nearby-cycle programme}
\label{sec:cclimax-chain-level-nearby-cycle-global-inversion}
\label{subsec:cclimax-global-inversion-climax-form}

\begin{theorem}[Nearby-cycle global inversion programme]
\label{thm:cclimax-global-inversion-all-admissible}
\ClaimStatusOpen
For every divisor \(H_n\) declared admissible, construct:
\begin{enumerate}[label=\textup{(\roman*)}]
\item the nearby-cycle CDG coalgebra and its monodromy filtration;
\item the local completed co-/contra-derived comparison;
\item an explicit extension class and its \(A_\infty\) null-homotopy;
\item derived Čech descent on all multiple intersections.
\end{enumerate}
Completion of these four tasks produces a global coderived inversion
theorem.  The admissible index set and each numerical wall coefficient
are outputs of the construction.
\end{theorem}

\begin{remark}[Evidence package for the nearby-cycle programme]
\label{rem:cclimax-four-verification-paths}
Local coefficient computations, Kashiwara--Malgrange specialization,
generic Koszul inversion, and automorphic divisor data test different
parts of Theorem~\ref{thm:cclimax-global-inversion-all-admissible}.
They share the recognition and comparison maps and therefore form one
coherent evidence package.
\end{remark}

\begin{remark}[Ambient-neutral statement of the programme]
\label{rem:cclimax-constitutional-ambient-neutral-theoremB}
The ambient-neutral conclusion is an isomorphism in a named coderived
category of a constructed completed CDG coalgebra.  The algebra,
coalgebra, completion, curvature, and descent topology are therefore
part of the theorem statement.
\end{remark}

\begin{remark}[Extension class]
\label{rem:cclimax-extension-obstruction-governing-equation}
\ClaimStatusOpen
A wall formula has the typed form
\[
 [d_{\mathrm{bar}},\phi^{(n)}]
 =
 \omega_n
 \quad\text{in the deformation complex of the counit}.
\]
Identifying \(\omega_n\) with a Fourier coefficient of an automorphic
lift requires a comparison map from that Fourier coefficient to the
same deformation complex.
\end{remark}

\section{Compatibility diagram for Theorems A, B, C, D, H}
\label{sec:climax-unified-five-theorem-identity}
\label{ssec:climax-unified-grand-synthesis}

\begin{theorem}[Five typed readouts of the chiral bar--cobar architecture]
\label{thm:unified-five-theorem-identity}
\ClaimStatusConditional
Let
\begin{equation}
\label{eq:unified-adjunction}
 \Omega_X^{\mathrm{ch}}
 :
 \ChirCoalg_X^{\mathrm{conil,comp}}
 \rightleftarrows
 \ChirAlg_X^{\mathrm{aug}}
 :
 B_X^{\mathrm{ch}},
 \qquad
 \Omega_X^{\mathrm{ch}}\dashv B_X^{\mathrm{ch}}
\end{equation}
be the adjunction on the ambient and completion locus of
Theorem~\ref{thm:A-infinity-2}.  Then:
\begin{enumerate}[label=\textup{(\Roman*)}]
\item Theorem~A concerns the unit and counit of
\eqref{eq:unified-adjunction}.
\item Theorem~B concerns strict and completed inversion.
\item Theorem~C concerns the Verdier-paired strict-flat
centre-local-system complex and its normalized scalar trace through
Theorem~D.
\item Theorem~D concerns modular pushforward of the genus-weighted
trace.
\item Theorem~H concerns the chiral Hochschild complex.
\end{enumerate}
Each item carries the hypothesis package of its pointer theorem.
Their common organizing diagram is
\eqref{eq:unified-adjunction}; their proofs remain the five native
arguments.
The computed and conditional scalar sets are
\begin{equation}
\label{eq:unified-kappa-plus-kappadual-finite-set}
 \{0,13,250/3\}_{\mathrm{Vol\,I}}
 \subset
 \{0,8,13,250/3\}_{H_{\mathsf B}}.
\end{equation}
\end{theorem}

\begin{proof}
The five clauses restate
Theorems~\ref{thm:theorem-a-platonic-pointer}--%
\ref{thm:theorem-h-platonic-pointer} with their type signatures.
Equation~\eqref{eq:unified-kappa-plus-kappadual-finite-set} is
Theorem~\ref{thm:theorem-c-platonic-pointer}.
\label{prf:unified-five-theorem-identity}
\end{proof}

\subsection{Five-archetype status table}
\label{ssec:climax-unified-five-archetypes}

\begin{proposition}[Five-archetype scalar status]
\label{prop:five-archetype-specialisation}
\ClaimStatusConditional
The archetype witnesses have the following current scalar and
shadow-depth data:
\[
\begin{array}{c|l|c|c|l}
\text{taxon}&\text{witness}&\text{scalar entry}&r_{\max}&\text{status}\\\hline
\mathsf G&\mathcal H_k&K^\kappa=0&2&\text{computed from partner formula}\\
\mathsf L&\widehat{\mathfrak g}_k&K^\kappa=0&3&\text{computed from partner formula}\\
\mathsf C&\beta\gamma_\lambda/bc_\lambda&K^\kappa=0&4&\text{computed from partner formula}\\
\mathsf M&\mathrm{Vir}_c&K^\kappa=13&\infty&\text{computed from partner formula}\\
\mathbf B&\widetilde H(K3,\mathbb Z)&K^{\mathrm{lat}}=8&5&
 \text{lattice theorem; chiral candidate under }H_{\mathsf B}
\end{array}
\]
The principal \(\mathcal W_3\) extension has computed value
\(K^\kappa=250/3\).  Bershadsky--Polyakov has proved \(K^c=50\) and
open \(K^\kappa\).  The crown normalization is
\begin{equation}
\label{eq:unified-crown-identity}
 K^{\mathrm{lat}}=8,\quad \hbar^2:=-1/8
 \quad\Longrightarrow\quad
 \hbar^2K^{\mathrm{lat}}=-1.
\end{equation}
The implication is scalar substitution; its quantum interpretation is
the \(H_{\mathrm{quantum}}\) obligation.
\end{proposition}

\begin{proof}
The first four rows and the principal-\(\mathcal W_3\) extension are
the family calculations of Chapter~\ref{ch:landscape-census}.  The
fifth row is the Mukai signature calculation.  The status column
records the comparison theorem required to pass from a lattice number
to a chiral conductor.
\end{proof}

\subsection{Verification obligations}
\label{ssec:climax-unified-three-verification-paths}

\begin{proposition}[Three verification obligations for the
compatibility diagram]
\label{prop:unified-identity-three-path-verification}
\ClaimStatusOpen
A proof identifying more than the organizing diagram requires:
\begin{enumerate}[label=\textup{(V\arabic*)}]
\item natural transformations connecting the five derived functors;
\item familywise specialization maps preserving the companion pairs;
\item a deformation-theoretic map from wall classes to the bar--cobar
counit.
\end{enumerate}
These are three separate obligations.  Their completion would promote
the compatibility diagram to a reconstruction theorem.
\end{proposition}

\begin{remark}[One architecture, five theorem packages]
\label{rem:unified-identity-one-theorem-five-shadows}
The adjunction~\eqref{eq:unified-adjunction} organizes the five
readouts.  A common organizing object and a common proof are distinct
mathematical assertions; the present chapter establishes the former
and records the natural transformations required for the latter.
\end{remark}

\begin{remark}[Typed cross-volume reading]
\label{rem:unified-identity-cross-volume-reading}
The Vol~II map \(\rho_\partial\) and the Vol~III composite
\(\mathrm{Sp}^{\mathrm{ch}}\circ\Phi^{\mathrm{FA}}\) enter through
their own packages \(H_{\mathrm{BL}}\) and \(H_\Phi\).  The Borcherds
weight \(\kappa_{\mathrm{BKM}}(\Delta_5)=5\), the Mukai lattice number
\(8\), and the chiral modular conductor are distinct scalar
invariants.
\end{remark}

\begin{remark}[Degree-three comparison]
\label{rem:climax-unified-chirhoch3-portrait}
The package $H_H(\cA;S)$ determines whether degree~$3$ occurs in the
curve-level chart complex.  The bounded Virasoro benchmark contains
degree~$3$.  CY-$d$, Gelfand--Fuks, Hall, and CoHA degree-three classes
belong to their respective complexes; a named chain map transports a
class to the chiral chart.
\end{remark}

\subsection{Reference-curve selection for the crown candidate}
\label{ssec:climax-crown-named-curve}

\begin{theorem}[Reference-curve selection obligation]
\label{thm:climax-crown-canonical-curve}
\ClaimStatusOpen
The specialization
\(\mathrm{Sp}^{\mathrm{ch}}_{\Sigma_2,C}\) takes the reference curve
\(C\) as part of its input.  For \(K3\times E\), the natural candidate
is the elliptic factor \(C=E\).  An elliptic K3 fibration
\(\pi\colon K3\to\mathbb P^1\) has base \(\mathbb P^1\) and, in the
generic case, twenty-four \(I_1\) discriminant points on that base.
The elliptic factor \(E\), the base \(\mathbb P^1\), and the nodal
fibres are three distinct geometric objects.

A theorem selecting a canonical chiral curve must construct
\(H_\Phi\), specify \((\Sigma_2,C)\), and prove independence under the
allowed choices.  The canonical curve therefore remains an output of
the comparison package.
\end{theorem}

\begin{corollary}[Conditional genus-one crown specialization]
\label{cor:climax-crown-genus-stratified}
\ClaimStatusConditional
Assume \(H_\Phi\) with \((\Sigma_2,C)=(K3,E)\) and assume
\(H_{\Delta_5}^{\mathrm{rec}}\).  Then the resulting chiral object is
defined on \(E\), and its bar--cobar and chiral Hochschild statements
are evaluated on \(\operatorname{Ran}(E)\).  The twenty-four
discriminant points of an elliptic K3 enter through an additional
correspondence from the base \(\mathbb P^1\).
\end{corollary}

\begin{remark}[Curve data and algebra data]
\label{rem:climax-crown-curve-richer}
Genus, marked points, and nodal fibres belong to the source geometry;
associator, coproduct, and completion belong to the algebra.  A
factorization theorem connecting them specifies the correspondence
between these two data sets.
\end{remark}

\begin{remark}[Named-module verification programme]
\label{rem:climax-named-module-scope-crown}
\ClaimStatusOpen
A module-level application of Theorem~B begins with constructed
comodules over the completed bar coalgebra and verifies the
co-/contra-derived comparison for each.  Proposed fundamental,
vacuum, adjoint, and imaginary-root modules become theorem witnesses
after their coactions, filtrations, and convergence are written
explicitly.
\end{remark}

\begin{remark}[Moduli selection principle]
\label{rem:climax-curve-selection-principle}
A rigidity statement for a marked curve requires a morphism to the
appropriate moduli stack and a calculation of its deformation
complex.  Group transitivity on a finite set of labels supplies a
symmetry of the marking; infinitesimal rigidity is the separate
vanishing statement \(H^1(C,T_C(-D))=0\) or its stack-theoretic
replacement.  This is the precise selection problem for a proposed
crown curve.
\end{remark}
